\raggedbottom
\chapter{Conclusión}

El desarrollo de BlockTrack ha permitido explorar y aplicar diversas tecnologías emergentes en el ámbito de la 
trazabilidad de logística de mercancías. A lo largo de este proyecto, se ha demostrado cómo la integración de blockchain,
puede ofrecer una solución robusta y segura para el seguimiento de productos a lo largo de la cadena de suministro.

Es cierto que, se ha aplicado en un caso de uso específico como es el transporte por carretera, pero los principios y 
tecnologías implementadas son fácilmente adaptables a otros modos de transporte y sectores industriales. La capacidad
de personalización y escalabilidad del sistema desarrollado abre la puerta a futuras expansiones y mejoras.

\section{Análisis del cumplimiento de los objetivos}
El objetivo general del proyecto es desarrollar una solución de trazabilidad logística utilizando tecnologías blockchain.
Como hemos visto, a lo largo del desarrollo de BlockTrack, en capítulos anteriores; como análisis de requisitos, diseño
de la arquitectura, implementación y pruebas, se ha cumplido con éxito este objetivo. La combinación de tecnologías usadas
para desarrollo de aplicaciones centralizadas y descentralizadas ha permitido crear una plataforma funcional y eficiente donde
el seguimiento de mercancías es transparente y seguro.

\section{Cumplimiento objetivos específicos}
\begin{itemize}
    \item \textbf{La aplicación deberá permitir la gestión completa de remesas: Cumplido.} La aplicación permite
gestionar remesas desde su creación hasta su entrega, incluyendo la actualización de estados y la visualización, además de operaciones básicas,
como edición de la misma.
    \item \textbf{ Uso de una red blockchain para registrar eventos clave en el ciclo de
vida de una remesa, garantizando la inmutabilidad y verificabilidad de
los datos: Cumplido.} A lo largo de este documento, se ha detallado y explicado cómo se ha implementado, estudiado y probado el uso de
blockchain para registrar eventos clave en el ciclo de vida de una remesa.
    \item \textbf{Diseñar e implementar contratos inteligentes que permitan registrar
y garantizar la trazabilidad de las remesas, asegurando su correcta
integración con la red blockchain y con la aplicación web: Cumplido.} Se han diseñado, implementado y probado contratos inteligentes para automatizar 
la gestión y trazabilidad de las remesas, asegurando su correcta integración con la red blockchain y la aplicación web.
    \item \textbf{Seguimiento de remesas en tiempo real, proporcionando a los usuarios
información actualizada sobre el estado de sus envíos: Cumplido.} Debido a las propiedades de descentralización de la red blockchain, cada vez que el usuario ejecuta una acción
sobre una remesa, esta desemboca en una acción automatizada por el smart contract que añade una nueva transacción a la cadena de blockchain. Debido a que cada usuario tiene una copia
actualizada de la cadena, el seguimiento en tiempo real está garantizado.
    \item \textbf{Implementar mecanismos de autenticación y autorización robustos para
     garantizar la seguridad y privacidad de los datos de los usuarios: Cumplido.} Se han implementado mecanismos de autentificación y autorización robustos utilizando
        tecnologías modernas como JWT y bcrypt, garantizando la seguridad y privacidad de los datos de los usuarios.
    \item \textbf{Desarrollar una interfaz de usuario intuitiva y accesible que facilite
la interacción con la aplicación y la visualización de la información
relevante: Cumplido.} La UI ha sido diseñada con un enfoque en la experiencia del usuario, utilizando frameworks 
modernos como React y TailwindCSS para crear una interfaz intuitiva y accesible. Siguiendo patrones como diseño móvil responsivo y
diseño basado en componentes, facilitándose así la coherencia y reutilización de elementos en toda la aplicación.
\end{itemize}

\section{Profundización en objetivos personales}
A nivel personal, este proyecto ha supuesto un reto significativo para mí, ya que me ha permitido
profundizar en áreas tecnológicas que eran nuevas para mí, como blockchain y contratos inteligentes. Ha sido
muy satisfactorio ver todo lo que he profundizado y aprendido a lo largo del desarrollo de BlockTrack sobre estas
tecnologías, y cómo he podido aplicarlas de manera práctica en un proyecto real. \\

Además, el proyecto me ha brindado la oportunidad de mejorar mis habilidades en desarrollo web, tanto en el lado
del frontend como del backend. He aprendido a trabajar con diversas herramientas y frameworks, lo que ha enriquecido
mi conjunto de habilidades y me ha preparado mejor para futuros desafíos en el campo del desarrollo de software. \\

En resumen, BlockTrack no solo ha cumplido con los objetivos planteados al inicio del proyecto, sino que también
ha sido una experiencia muy gratificante a nivel personal y profesional aunque en determinados momentos haya
supuesto un desafío.

\section{Trabajos futuros}
Aunque BlockTrack ha logrado cumplir con sus objetivos iniciales, ha conseguido probar que la tecnología blockchain
puede ser una solución efectiva para la gestión de remesas. Sin embargo, existen varias áreas en las que se podría
mejorar y expandir el proyecto en el futuro:

\begin{itemize}
    \item \textbf{Integración con IoT:} La incorporación de dispositivos IoT para el seguimiento en tiempo real de las
    condiciones de las mercancías (temperatura, humedad, ubicación GPS) podría mejorar significativamente la trazabilidad y
    seguridad de las remesas.
    \item \textbf{Escalabilidad:} Investigar y aplicar soluciones de escalabilidad para blockchain, como redes de segunda capa,
    podría mejorar el rendimiento y reducir los costos de transacción a medida que la plataforma crece.
    \item \textbf{Gestión de aduanas:} Implementar funcionalidades específicas para la gestión de aduanas y regulaciones internacionales
   podría ampliar el alcance a nivel global. La blockchain podría utilizarse para almacenar y verificar documentos aduaneros de manera segura.
    \item \textbf{Análisis de datos:} Desarrollar herramientas de análisis de datos para proporcionar a los usuarios información
    valiosa sobre sus remesas, como patrones de envío, tiempos de entrega y posibles áreas de mejora en la cadena de suministro.
    \item \textbf{Mejora de la interfaz de usuario:} Continuar mejorando la experiencia del usuario mediante la incorporación
    de nuevas funcionalidades y optimizaciones basadas en el feedback de los usuarios.
    \item \textbf{Soporte para distintos tipos de transporte:} Ampliar la plataforma para soportar otros modos de transporte,
    como marítimo, aéreo y ferroviario, para ofrecer una solución más completa de trazabilidad logística.
\end{itemize}

\section{Planificación final ajustada}
A lo largo del desarrollo de BlockTrack, la planificación inicial ha sufrido varios ajustes debido a diversos factores,
como la complejidad técnica de ciertas implementaciones y la necesidad de dedicar más tiempo a la investigación de blockchain, Ethereum y smart contracts.\\

También influyó que en medio del desarrollo del proyecto, comencé a trabajar en una empresa del sector tecnológico, lo que
afectó a mi disponibilidad de tiempo para dedicar al proyecto. \\

Ajustes en los requisitos iniciales también llevaron a cambios en la planificación, ya que algunas funcionalidades se añadieron o modificaron
durante el desarrollo como por ejemplo por falta de conocimiento previo, no tuve en cuenta que toda la información referente al envío quedaría registrada en la blockchain por lo que
no sería necesario almacenar ciertos datos en la base de datos centralizada. Llevando a cambios en el diagrama E-R, paso a tablas,...\\

La planificación definitiva es la siguiente:
\begin{table}[H]
\centering
\renewcommand{\arraystretch}{1.3}
\setlength{\tabcolsep}{8pt}
\begin{tabular}{|p{8cm}|c|}
\hline
\textbf{Fase} & \textbf{Duración (días)} \\ \hline
Investigación sobre el problema & 14 \\ \hline
Análisis de requisitos & 20 \\ \hline
Investigación sobre las tecnologías más eficaces para el desarrollo & 10 \\ \hline
Escritura de la memoria técnica & 60 \\ \hline
Desarrollo MVP & 88 \\ \hline
Investigación soluciones planificación de rutas & 10 \\ \hline
Investigación tecnología blockchain & 30 \\ \hline
Investigación Ethereum, smart contracts y como conectarlo con backend y frontend & 20 \\ \hline
Implementación de smart contract y ethereum & 20 \\ \hline
Corrección de bug & 17 \\ \hline
\textbf{TOTAL} & \textbf{289 DÍAS} \\ \hline
\end{tabular}
\caption{Planificación temporal del proyecto}
\label{tab:planificacion}
\end{table}

\section{Reflexión final}
Para terminar, me gustaría destacar que este proyecto ha sido una experiencia de aprendizaje invaluable. Me ha permitido
adentrarme en el mundo de las tecnologías blockchain y su aplicación práctica en la logística, 
lo cual considero que es un campo con un gran potencial de crecimiento en el futuro y con todavía frentes abiertos para resolver.\\

Considero que este TFG da como resultado una demostración sólida de cómo las tecnologías emergentes pueden integrarse en
soluciones prácticas para problemas del mundo real, fascinándome cómo de fácil permite blockchain mejorar la trazabilidad y seguridad en
la gestión de remesas. Estoy entusiasmado por las posibilidades futuras que estas tecnologías ofrecen y espero seguir explorándolas en mi carrera profesional. 

